\section{Conclusion}
\label{sec:conclusion}

We presented Strided DLLM, a diffusion language model that bridges autoregressive and parallel decoding by respecting the pretrained AR model's structure at every level: causal attention, next-token prediction, and serving infrastructure. The unified verify-and-decode training objective equips the model with both a decode distribution $q$ (from masked positions) and a verify distribution $p$ (from clean positions), enabling Introspective Strided Decoding---a principled mechanism that adaptively selects the effective stride per step while provably preserving the base AR distribution.

On the systems side, we demonstrated that high-performance DLLM serving does not require a fundamentally different infrastructure stack. By mapping ISD to SGLang's native extend mode and introducing a stationary-batch decode loop, CUDA graph capture for multi-token extend, fused introspective verification, and acceptance-maximizing proposals, Strided DLLM-8B achieves 2.1$\times$ the per-request throughput of Qwen3-8B AR at $C{=}1$ and sustains advantages through $C{=}64$ on a single H100---outperforming both the AR baseline and EAGLE-3 across all concurrency levels. For applications requiring exact AR-quality output, segment-gated LoRA with cuBLAS kernel replacement and multi-stream overlap delivers bit-for-bit lossless acceleration at 76\% of the non-LoRA throughput.

Our analysis identifies the serial dependency chain within each ISD step as the primary remaining bottleneck, limiting multi-GPU scaling to 1.5$\times$ (TP=1$\to$4) compared to AR's 2.15$\times$. We outlined three directions to address this: speculative metadata pre-computation, hybrid decode--extend attention, and pipelined verification. We will release model weights, training code, and the SGLang integration to facilitate further research and deployment.
